\chapter{The Examination of the Amber Gate}

A sample scenario for Game Masters. This is the practical examination I give to my advanced students. It is not a combat. It is a test of grammar, judgement, and the willingness to accept a receipt.

\section{Setup}

The candidate stands in a sealed chamber. The chamber has one door, currently locked. On a table in the center of the room rests a single candle---unlit. On the wall opposite the door hangs a mirror. The mirror reflects the candidate's face, but the reflection is slightly delayed.

The candidate is told: \emph{``Leave the room. You may not touch the door. You may not break the mirror. The candle is your only tool. You have until the candle burns out.''}

The candle is not magical. It will burn for exactly one scene.

\section{The Challenge}

The door is locked with a mundane lock. The candidate cannot touch the door. The mirror is a \emph{witness}. The Weave will remember anything the mirror sees.

The candidate must use the candle and the mirror to unlock the door from a distance. There are multiple solutions.

\subsection{Solution One: Veil and Leap}

The candidate casts [VEIL] on the candle's flame, making it invisible to the mirror. The mirror does not witness the next action. The candidate then casts [LEAP] to teleport their own hand through the door (line of sight required---but the door is wood; they cannot see through it). This fails unless they also cast [SCRY] to see the far side. Three tags. DV 5. Risky.

\subsection{Solution Two: Memory and Mirror}

The candidate lights the candle. The mirror reflects the flame. The candidate casts [MEMORY] on the mirror, instructing it to \emph{forget} that the door is locked. The Weave accepts this if the candidate has touched the lock previously (so the mirror has a memory of it unlocked). DV 4. Requires preparation.

\subsection{Solution Three: The Direct Approach}

The candidate casts [TRANSFORM] on the lock, turning it into wax. The lock melts. The door opens. The mirror witnesses everything. The Weave charges a receipt for the transformation. The candidate accepts it. DV 5 (dangerous). Simple. Expensive.

\subsection{Solution Four: The Student Who Read the Assignment}

The candidate notices that the instruction says \emph{``You may not touch the door. You may not break the mirror. The candle is your only tool.''} It does not say the candidate cannot ask for help. The candidate knocks. The examiner opens the door from the outside.

The candidate passes with highest honors.

\section{Scoring}

\begin{itemize}
    \item \textbf{Clean success (any solution)}: Pass.
    \item \textbf{Partial success}: Pass with a note to review tag combinations.
    \item \textbf{Miss, minor receipt}: Retake in one month.
    \item \textbf{Miss, major receipt}: Retake after remedial Grammar.
    \item \textbf{Catastrophic failure}: The examination hall is on fire. The candidate is awarded a passing grade for spectacle.
\end{itemize}

\section{GM Notes}

This examination tests creative problem- solving, not combat. Encourage players to propose unusual tag combinations. Reward clever interpretations of the instructions. The Weave respects cleverness almost as much as it respects precision.

If the candidate sets the mirror on fire, the mirror will remember. The Weave will include that memory in its receipt.

\textit{--- Lysandra}
